% arXiv-compatible preprint of paper/workflow/parallel-dispatch-breakeven-point.
% Source-of-truth for the content remains the PAPER.md frontmatter + body in the
% parent directory. This LaTeX file is a hand-rendered preprint counterpart that
% mirrors the IMRaD body for venues that need a PDF.
%
% Build:
%   pdflatex paper.tex
%   bibtex paper
%   pdflatex paper.tex
%   pdflatex paper.tex
%
% Or invoke the helper:  bash build.sh
%
% Class choice: standard `article` keeps the source portable to arXiv, OpenReview,
% and most workshop templates. Switch to `acmart` / `IEEEtran` / a venue-specific
% class only when an actual venue is targeted.

\documentclass[10pt,letterpaper]{article}

\usepackage[margin=1in]{geometry}
\usepackage[utf8]{inputenc}
\usepackage{microtype}
\usepackage{booktabs}
\usepackage{array}
\usepackage{hyperref}
\usepackage{xcolor}
\usepackage{verbatim}
\usepackage{amsmath}

\hypersetup{
  colorlinks=true,
  linkcolor=blue!50!black,
  citecolor=blue!50!black,
  urlcolor=blue!70!black,
}

\title{When Does Parallel Subagent Dispatch Stop Paying Off?\\
\large A useful-output cost gate for bulk-modification workloads in LLM agent pipelines}

\author{kjuhwa\thanks{kjuhwa@nkia.co.kr; \texttt{kjuhwa/skills-hub}}}

\date{April 24, 2026 \quad (preprint v1, IMRaD body finalized 2026-04-25)}

\begin{document}
\maketitle

\begin{abstract}
A bulk-modification task dispatched to $N$ parallel LLM subagents pays a fixed
startup overhead per agent plus a verification cost per file every agent must
read before deciding whether to act. When prior-work coverage is high, the
dispatch can return zero useful changes from most agents while still paying
the full overhead. The original draft of this work proposed a 70\,\%
prior-coverage threshold above which parallel dispatch is net negative.
A cost-model harness applied to five corpus domains in
\texttt{kjuhwa/skills-hub} (range $36.8\,\%$--$100\,\%$ coverage) refutes
the 70\,\% claim as a universal constant and identifies a more robust gate:
\textbf{absolute useful-output count}, not coverage fraction. The revised
criterion is that parallel dispatch is unjustified when fewer than
$\sim 5$ files actually need real work, regardless of coverage. The
directional claim (very high coverage trends toward non-parallel) survives;
the specific threshold does not. We summarize the cost model, the workload
selection, and the implications for tool design when the catalog of
``how to parallelize'' skills lacks a corresponding ``when not to'' gate.
\end{abstract}

\section{Introduction}

A bulk-modification task dispatched to $N$ parallel subagents pays a fixed
startup overhead per agent plus a verification cost per file each agent must
\emph{Read} before deciding whether to act. When prior-work coverage is high,
the dispatch can return zero useful changes from most agents while still
paying the full overhead. Existing parallelism guidance focuses on
\emph{how} to dispatch but not on \emph{when not to}.

This paper sets out to identify the gate criterion: a single rule a
pre-flight check can apply to refuse a dispatch when waste dominates. The
original draft proposed a 70\,\% coverage threshold; the experiment in this
paper tests that claim against measured corpus data.

\subsection{Background}

In the host catalog, three procedure-level entries describe parallel
dispatch:
\begin{itemize}
\item \texttt{workflow/parallel-bulk-annotation} --- canonical \emph{how-to}
  for annotating many files across parallel agents.
\item \texttt{workflow/bucket-parallel-java-annotation-dispatch} --- a
  language-specific variant using bucket partitioning.
\item \texttt{ai/ai-subagent-scope-narrowing} --- the adjacent concept of
  narrowing what gets dispatched.
\end{itemize}

And one pitfall, authored mid-session when the pattern broke:
\begin{itemize}
\item \texttt{agent-orchestration/grep-existing-annotations-before-parallel-subagent-dispatch}
  --- four parallel agents dispatched to annotate Java DTO files; three
  returned ``zero changes --- already annotated''; three agent rounds
  wasted because prior coverage was not checked first.
\end{itemize}

The pitfall is the counter-evidence for this work. The skills tell you
\emph{how} to parallelize; the pitfall tells you \emph{when it stops paying
off}. The two pieces are split across kinds and across files; nothing
enforces that a user reading the skill also encounters the pitfall.

\subsection{Premise revision (2026-04-24)}

The original premise stated a \emph{70\,\% coverage threshold}. The
experiment in \S\ref{sec:methods}--\S\ref{sec:results} tested that claim
against five measured corpus domains and found the 70\,\% number was not a
robust constant --- it shifts with the $\alpha/w$ cost ratio. The meaningful
gate turned out to be \textbf{absolute useful-output count}, not coverage
fraction.

\subsection{Prior art}

External references are not yet cited here. Relevant directions include
parallel-algorithm amortization literature on fork-join break-even, queueing
theory cost models adapted for agent-as-worker, LLM token-cost analysis in
developer tooling, and pre-flight gate patterns in other agent frameworks
(LangGraph, AutoGen, Crew). Filling these in would ground the threshold
claim in prior art rather than treating it as intuition; this paper makes
no claim of completeness on that front.

\section{Methods}
\label{sec:methods}

\subsection{Cost model}

Total dispatch cost is approximately
\begin{equation}
C_{\text{total}} = N\alpha + Nv\frac{T}{N} + w \cdot \mathrm{useful},
\end{equation}
where $\alpha$ is per-agent startup overhead, $v$ is per-file
verification cost, $w$ is per-file real-work cost, $T$ is total file
count split across $N$ agents, and $\mathrm{useful} = (1 - c) \cdot T$
is the count of files that actually need work given prior coverage
$c \in [0, 1]$.

Wall-clock for parallel dispatch is approximately
\begin{equation}
W_{\parallel} \approx \alpha + v \frac{T}{N} + w \frac{\mathrm{useful}}{N},
\end{equation}
while the serial baseline is
\begin{equation}
W_{\text{serial}} \approx \alpha + vT + w \cdot \mathrm{useful}.
\end{equation}
The crossover is where $C_{\text{total}}$ stops dominating
$C_{\text{serial}}$, driven by the $\alpha/w$ ratio and the absolute
$\mathrm{useful}$ count, \emph{not} by $c$ alone. Coverage fraction is a
proxy that fails when total $T$ varies between domains.

\subsection{Pre-flight probe}

The probe is one second of shell:
\begin{verbatim}
total=$(find <scope> -name "*.java" | wc -l)
hit=$(grep -rln "@Schema" <scope> | wc -l)
coverage=$((hit * 100 / total))
useful=$((total - hit))
\end{verbatim}
The decision can fire on \texttt{useful} (absolute count) or
\texttt{coverage} (fraction). The experiment tests which of the two is the
more robust gate criterion in practice.

\subsection{Workload selection}
\label{sec:workloads}

Five domains within the host catalog were measured via grep against the
remote cache (Table~\ref{tab:workloads}).

\begin{table}[h]
\centering
\caption{Measured corpus domains (April 2026 snapshot).}
\label{tab:workloads}
\begin{tabular}{lllrr}
\toprule
ID & Domain & Marker & Total & Hit \\
\midrule
A & skills with \texttt{triggers:} populated & YAML key non-empty & 468 & 239 \\
B & skills with \texttt{content.md} sibling & file exists & 468 & 172 \\
C & knowledge with \texttt{description:} & YAML key non-empty & 351 & 193 \\
D & knowledge with \texttt{tags:} & YAML key present & 351 & 351 \\
E & skills with stable version $\geq 1.0$ & semver gate & 468 & 376 \\
\bottomrule
\end{tabular}
\end{table}

\subsection{Cost-model harness}

A small classifier (\texttt{example/workflow/coverage-gate-benchmark} in
the host catalog) consumes the tuple $(N, \alpha, v, w, \mathrm{useful}, T)$
and emits one of \texttt{PARALLEL}, \texttt{BORDERLINE}, \texttt{SAMPLING},
\texttt{NO DISPATCH}. Default weights for this run:
$\alpha = 30$\,s, $v = 5$\,s, $w = 60$\,s, $N = 4$. The classifier output is
compared per-domain against the original 70\,\%-coverage rule.

\section{Results}
\label{sec:results}

The 70\,\% coverage threshold as a \emph{universal} constant is not
supported. Under the default weights, parallel wall-clock savings dominate
up to $\sim 90\,\%$ coverage for any domain with non-zero useful output.
The crossover shifts with the $\alpha/w$ ratio; it is not a robust
constant.

\begin{table}[h]
\centering
\caption{Per-domain classification under cost model vs original 70\,\% rule.}
\label{tab:results}
\begin{tabular}{lrrll}
\toprule
ID & Coverage $c$ & Useful & Cost-model verdict & 70\,\%-rule verdict \\
\midrule
A & 51.1\,\% & 229 & \textsc{PARALLEL} & \textsc{PARALLEL} \\
B & 36.8\,\% & 296 & \textsc{PARALLEL} & \textsc{PARALLEL} \\
C & 55.0\,\% & 158 & \textsc{PARALLEL} & \textsc{PARALLEL} \\
D & 100.0\,\% & 0 & \textsc{NO DISPATCH} & \textsc{NO DISPATCH} \\
E & 80.3\,\% & 92 & \textbf{\textsc{PARALLEL}} & \textsc{SAMPLING} (mismatch) \\
\bottomrule
\end{tabular}
\end{table}

Domain D (zero useful output, full coverage) is correctly rejected by both
rules. Domain E exposes the discrepancy: at 80.3\,\% coverage, the original
rule says ``switch to sampling,'' but the cost model says ``still
parallelize'' because $92 \times w = 92$ minutes of real work dwarfs
$N \times \alpha = 2$ minutes of agent startup. Coverage fraction alone
misclassifies.

The discovered gate is not coverage percentage but
\textbf{useful-output absolute count}: when fewer than $\sim 5$ files
actually need work, parallel dispatch is pure waste regardless of $c$. The
directional claim (very high coverage trends toward non-parallel) holds;
the specific 70\,\% number was a single-session observation, not a robust
threshold. We classify the verdict as \emph{partial support} and rewrite
the premise to match.

\section{Discussion}
\label{sec:discussion}

\subsection{Why coverage fraction misled}

Coverage is a ratio, not an absolute. A domain with 92 useful files at
80\,\% coverage has the same parallel justification as a domain with 92
useful files at 30\,\% coverage --- the parallel work scales with the
absolute, the agent startup is paid once. The original premise inherited
an intuition from one pitfall session (where useful output was small
because total was small) and over-generalized to fraction.

\subsection{Asymmetric discoverability}

``How to parallelize'' is a teachable procedure that composes well as a
catalog skill. ``When not to parallelize'' is a judgment call requiring
a cost-curve mental model the author has to encode. The host catalog has
many \emph{how-to} skills and one pitfall; these are asymmetric in
discoverability --- how-to skills get triggered by keyword matching
(``bulk annotation''), the pitfall gets found only if the user searches
for ``wasted agent parallel.'' The proposed remedy is a dedicated gate
skill whose entire job is the decision, so the gate is found first.

\subsection{Failure modes are silent}

A parallel dispatch that wastes tokens fails silently --- all $N$ agents
return success, wall-clock looks similar to the productive case, the only
signal is the token bill. Silent failure accumulates across sessions. The
gate converts silent waste into a visible decision: \emph{``coverage
92\,\%, useful 3 --- parallel not justified, switch to sampling.''} The
user can override, but the override is explicit and auditable.

\subsection{Limitations}

\begin{itemize}
\item \textbf{$n=1$ original evidence.} The 70\,\% threshold rested on one
  observed session. Even after the experiment runs, the data points are
  five hub-internal domains; cross-organization replication is needed
  to claim full generality.
\item \textbf{Cost weights are illustrative.} $\alpha$, $v$, $w$ are
  plausible values from agent-runner observation, not measured per
  deployment. A team with much faster startup or much slower
  verification will see a different crossover.
\item \textbf{Gate skill is unbuilt.} The proposed builds are
  POC-scoped; the paper is a call for work, not a claim of completed
  work.
\item \textbf{Counter-argument unexplored.} In deployments where token
  costs are negligible (local models, fixed-cost inference), the gate's
  value drops. The assumption that tokens cost enough to justify the
  gate may not hold universally.
\end{itemize}

\subsection{Future work}

\begin{enumerate}
\item Re-run the harness against domains in 2--3 other repositories to
  test cross-corpus stability of the useful-output gate.
\item Decide whether the coverage probe should be automatic (a hook on
  any \texttt{/hub-make --parallel} invocation) or opt-in via the gate
  skill. Automatic risks false positives; opt-in risks being forgotten.
\item Decide whether existing parallel skills should be \emph{retracted}
  and replaced with gate-mediated variants, or just annotated with a
  ``Phase 0 --- pre-flight'' section. The annotation is pragmatic; the
  retraction is more honest about the asymmetry.
\end{enumerate}

\section*{Provenance and source of truth}

This is a preprint counterpart of the source-of-truth paper at\\
\texttt{paper/workflow/parallel-dispatch-breakeven-point/PAPER.md} in the
\texttt{kjuhwa/skills-hub} repository. The PAPER.md frontmatter carries
\texttt{premise.if/then}, \texttt{examines[]}, \texttt{experiments[]} (with
\texttt{result}, \texttt{supports\_premise}, \texttt{observed\_at}), and
\texttt{outcomes[]} entries (\texttt{produced\_knowledge:
decision/parallel-dispatch-useful-output-gate};
\texttt{produced\_example: workflow/coverage-gate-benchmark}).
The body adopts the IMRaD structure documented in
\texttt{docs/rfc/paper-schema-draft.md} \S 5. The status is
\texttt{implemented}: experiments[0] completed with
\texttt{supports\_premise: partial} and the corpus changes above were
shipped.

\paragraph{Schema and tooling.}
The paper schema, the auto-injected references block, the IMRaD
compliance audit, and the citations index that powers the host catalog's
discovery flow are all open-source under the same repository. See the
\texttt{bootstrap/tools/} directory for the audit and indexer pipeline.

\bibliographystyle{plain}
\bibliography{references}

\end{document}
